\subsection{Листинг численного метода}

Для реализации блок-схем воспользуемся языком программирования C++.

\begin{lstlisting}[caption=Файл \texttt{differential.hpp}]
#ifndef DIFFERENTIAL_H
#define DIFFERENTIAL_H

#include "preamble.h"
#include <cstdlib>

std::size_t ITERATIONS_MAX = 1e6;

struct EulerSolution {
  double answer;
  std::size_t n_grid;
};

double calcForGrid(int f, double y_a, double a, double b, std::size_t n_grid) {
  double h = (b - a) / n_grid;
  std::function<double(double,double)> func = get_function(f);
  double y_b = y_a;
  double x = a;
  for (std::size_t i = 1; i <= n_grid; i++) {
    y_b += h * func(x, y_b);
    x += h;
  }
  return y_b;
}

EulerSolution solveByEuler(int f, double epsilon, double a, double y_a, double b) {
  std::size_t n_grid = 2;
  double y_b_h = calcForGrid(f, y_a, a, b, n_grid);
  while (n_grid < ITERATIONS_MAX) {
    n_grid *= 2;
    double y_b_h_2 = calcForGrid(f, y_a, a, b, n_grid);
    if (std::abs(y_b_h - y_b_h_2) <= epsilon) {
      return { y_b_h_2, n_grid };
    }
    y_b_h = y_b_h_2;
  }
  return { y_b_h, n_grid };
}

#endif
\end{lstlisting}

\begin{lstlisting}[caption=Файл \texttt{interpolate.hpp}]
#ifndef INTERPOLATE_H
#define INTERPOLATE_H

#include <vector>
#include <tuple>
#include <functional>
#include <algorithm>
#include <cmath>

struct InterpolationResult {
  std::vector<std::tuple<double, double>> points;
  std::size_t chebyshev_k;
};

std::vector<double> GetChebyshevRoots(double a, double b, std::size_t n_size) {
  std::vector<double> roots(n_size);
  const double pi = std::acos(-1.0);

  for (std::size_t i = 0; i < n_size; i++) {
    double x_standard = std::cos((2.0 * i + 1.0) / (2.0 * n_size) * pi);
    roots[i] = 0.5 * (a + b) + 0.5 * (b - a) * x_standard;
  }

  return roots;
}

InterpolationResult InterpolateLagrange(double a, double b, std::function<double(double)> func, double eps, std::size_t n_lagrange_max, std::size_t dots_num) {
  std::vector<std::tuple<double, double>> lagrange_values(dots_num + 1);
  std::vector<double> prev_values(dots_num + 1, 0.0);
  std::vector<double> curr_values(dots_num + 1, 0.0);

  auto eval_lagrange = [](double x, const std::vector<double>& nodes, const std::vector<double>& values) {
    double total_sum = 0.0;
    std::size_t n = nodes.size();
    for (std::size_t i = 0; i < n; ++i) {
      double term = values[i];
      for (std::size_t j = 0; j < n; ++j) {
        if (j != i) {
          term *= (x - nodes[j]) / (nodes[i] - nodes[j]);
        }
      }
      total_sum += term;
    }
    return total_sum;
  };

  std::size_t final_k = 0;
  for (std::size_t k = 5; k <= n_lagrange_max; ++k) {
    final_k = k;
    auto nodes = GetChebyshevRoots(a, b, k);
    std::vector<double> values(k);
    for (std::size_t j = 0; j < k; ++j) {
      values[j] = func(nodes[j]);
    }

    for (std::size_t i = 0; i <= dots_num; ++i) {
      double x = (i == 0) ? a : ((i == dots_num) ? b : a + i * (b - a) / dots_num);
      curr_values[i] = eval_lagrange(x, nodes, values);
    }

    if (k > 1) {
      double max_diff = 0.0;
      for (std::size_t i = 0; i <= dots_num; ++i) {
        max_diff = std::max(max_diff, std::abs(curr_values[i] - prev_values[i]));
      }
      if (max_diff <= eps) {
        break;
      }
    }
    prev_values = curr_values;
  }

  for (std::size_t i = 0; i <= dots_num; ++i) {
    double x = (i == 0) ? a : ((i == dots_num) ? b : a + i * (b - a) / dots_num);
    lagrange_values[i] = std::make_tuple(x, curr_values[i]);
  }

  return {lagrange_values, final_k};
}

#endif
\end{lstlisting}

\begin{lstlisting}[caption=Файл \texttt{main.cpp}]
#include <iostream>
#include <vector>
#include <cmath>
#include <iomanip>
#include <fstream>
#include <string>
#include <filesystem>

#include "lib/differential.hpp"
#include "lib/interpolate.hpp"

struct TestCase {
  int func_number;
  int case_index;
  double a;
  double y_a;
  double b;
  double epsilon;
};

double getExactResult(int f, double a, double y_a, double b) {
  switch (f) {
    case 1:
      return y_a + std::cos(a) - std::cos(b);
    case 2:
      return y_a * std::exp((b * b - a * a) / 4.0);
    case 3: {
      double C = (y_a * y_a - 2.0 * a - 1.0) * std::exp(-2.0 * a);
      double expr = 2.0 * b + 1.0 + C * std::exp(2.0 * b);
      return (expr >= 0.0) ? std::sqrt(expr) : std::nan("");
    }
    case 4:
      return (y_a + a + 1.0) * std::exp(b - a) - b - 1.0;
    case 5: {
      double denom = (1.0 / y_a) - (b - a);
      return (std::abs(denom) > 1e-9) ? (1.0 / denom) : std::nan("");
    }
    default:
      return 0.0;
  }
}

void GenerateTrajectoryData(const TestCase& t) {
  namespace fs = std::filesystem;

  fs::path dir_path = "output";
  fs::create_directories(dir_path);

  std::string prefix = "f" + std::to_string(t.func_number) + "_case" + std::to_string(t.case_index);
  fs::path euler_path = dir_path / (prefix + "_euler.txt");
  fs::path exact_path = dir_path / (prefix + "_exact.txt");

  std::ofstream euler_file(euler_path);
  std::ofstream exact_file(exact_path);

  if (!euler_file.is_open() || !exact_file.is_open()) return;

  std::size_t samples = 100;
  std::size_t euler_iterations = 0;

  auto eval_euler = [t, &euler_iterations](double x) {
    EulerSolution solution = solveByEuler(t.func_number, t.epsilon, t.a, t.y_a, x);
    euler_iterations = std::max(euler_iterations, solution.n_grid);
    return solution.answer;
  };
  auto eval_exact = [t](double x) {
    return getExactResult(t.func_number, t.a, t.y_a, x);
  };

  InterpolationResult euler_result = InterpolateLagrange(t.a, t.b, eval_euler, t.epsilon, samples, samples);
  InterpolationResult exact_result = InterpolateLagrange(t.a, t.b, eval_exact, t.epsilon, samples, samples);
  euler_file << euler_result.chebyshev_k << " " << t.a << " " << t.b << " " << t.epsilon << " " << t.y_a << " " << euler_iterations << '\n';
  exact_file << exact_result.chebyshev_k << '\n';
  for (std::tuple<double, double> point : euler_result.points) {
    euler_file << std::get<0>(point) << " " << std::get<1>(point) << "\n";
  }
  for (std::tuple<double, double> point : exact_result.points) {
    exact_file << std::get<0>(point) << " " << std::get<1>(point) << "\n";
  }
}

int main() {
  std::cout << std::fixed << std::setprecision(6);

  std::vector<TestCase> tests = {
    {1, 1, 0.0, 0.0, 3.141592, 1e-4}, 
    {1, 2, 0.0, 1.0, 6.283185, 1e-4}, 
    {2, 1, 0.0, 1.0, 2.0, 1e-4}, 
    {2, 2, 0.0, 0.2, 7.2, 1e-4}, 
    {3, 1, 0.0, 2.0, 1.5, 1e-4}, 
    {3, 2, 0.0, 1.0, 0.8, 1e-4}, 
    {4, 1, 0.0, 0.0, 1.0, 1e-4}, 
    {4, 2, 0.0, -1.0, 2.0, 1e-4}, 
    {5, 1, 0.0, 0.5, 1.5, 1e-4}, 
    {5, 2, 0.0, 0.8, 1.1, 1e-4},
  };

  for (const TestCase& t : tests) {
    GenerateTrajectoryData(t);
    std::cout << "Trajectory data written to disk.\n";
  }

  return 0;
}
\end{lstlisting}